\emph{Setting.} $V = F \otimes \Lambda^{\bullet}W$ with $F = \Lambda^{\bullet}D$ is the composite space of Theorem \claimref{synthesis.joint-clifford-clause-is-a-klein-convention-on-the-composite-space}, with its Fock inner product (Theorem \claimref{std-algebra.fock-representation}); there the Klein-twisted slot operators $\tilde a_{1}, \tilde a_{2}$ and the pair operators $b_{1}, b_{2}, b_{3}$ satisfy the canonical anticommutation relations of one system of five modes, with adjoints the Fock adjoints (its (i), (ii)), and the untwisted $a_{i}$ commute with the $b_{p}$. In the occupation basis of $\Lambda^{\bullet}\C^{5}$ these are the Jordan-Wigner operators of five modes (slots first), and the check works there: $\Lambda^{\bullet}\C^{5}$ has the rational occupation basis and every Majorana combination below is a signed permutation matrix in it, so all statements are about real matrices and need no imaginary unit.

\emph{(i).} Write $\gamma = \alpha + \alpha^{*}$, $\gamma' = \alpha - \alpha^{*}$ for each mode, so $\alpha = \tfrac12(\gamma + \gamma')$, $\alpha^{*} = \tfrac12(\gamma - \gamma')$. Then $\gamma^{2} = \alpha^{2} + \alpha^{*2} + \{\alpha, \alpha^{*}\}$, $\gamma'^{2} = \alpha^{2} + \alpha^{*2} - \{\alpha, \alpha^{*}\}$ and $\{\gamma, \gamma'\} = 2(\alpha^{2} - \alpha^{*2})$, and for two modes $\alpha$, $\beta$ the four anticommutators of $\gamma_{\alpha}, \gamma'_{\alpha}$ with $\gamma_{\beta}, \gamma'_{\beta}$ are the images of $\{\alpha, \beta\}$, $\{\alpha, \beta^{*}\}$, $\{\alpha^{*}, \beta\}$, $\{\alpha^{*}, \beta^{*}\}$ under an invertible linear map (the change of basis above, applied to each factor). So the CAR ($\alpha^{2} = 0$, $\{\alpha, \alpha^{*}\} = 1$, cross anticommutators zero, and the relations with $\alpha^{*}$ obtained by taking adjoints) hold exactly when $\gamma^{2} = 1$, $\gamma'^{2} = -1$ for every mode and every two distinct Majoranas anticommute. On $V$, $-1$ is central with square $1$. If the CAR hold, the group $G$ generated by the ten $\gamma$'s and $-1$ consists of the $\pm\gamma_{S}$, $\gamma_{S}$ the increasing products; a nonempty $\gamma_{S}$ is traceless on $V$ (it anticommutes with an invertible $\gamma_{j}$, $j \notin S$ when $|S|$ is odd, which exists since $|S| < 10$, and $j \in S$ when $|S|$ is even, so $\operatorname{tr}\gamma_{S} = \operatorname{tr}(\gamma_{j}^{-1}\gamma_{S}\gamma_{j}) = -\operatorname{tr}\gamma_{S}$), while $\pm1$ has trace $\pm32$, so no nonempty $\gamma_{S}$ is $\pm1$ and the $\gamma_{S}$ are pairwise trace-orthogonal ($\gamma_{S}^{-1}\gamma_{T} = \pm\gamma_{S \triangle T}$). Hence $G/\{\pm1\}$ is elementary abelian of order $2^{10}$ with basis the images of the $\gamma$'s, distinct directions satisfy $\gamma_{i}\gamma_{j} = (-1)\gamma_{j}\gamma_{i}$, and $-1 \ne 1$: (F2) holds for the ten directions with $c = -1$ and $e_{i}^{2} = c^{s_{i}}$, $s = (0, 1, 0, 1, \dots)$, and the representation of $\Z[G]$ on $V$ sends $c$ to $-1$, so it factors through $A = \Z[G]/(c + 1)$ (F3). By Theorem \claimref{synthesis.signed-directions-and-the-integer-clifford-algebra}(v), (vii), $A = \Cl_{\Z}(5,5)$, and it acts faithfully on $V$ by trace-orthogonality of the blades. Conversely, if the ten operators with $-1$ form a group of signed directions with these squares, the identities above give the CAR for $\tfrac12(\gamma + \gamma')$ and $\tfrac12(\gamma - \gamma')$; the factor $\tfrac12$ needs the rationals, which the coefficients of the operators on $V$ already are. For the given modes the CAR hold (Theorem \claimref{synthesis.joint-clifford-clause-is-a-klein-convention-on-the-composite-space}(ii), and Theorem \claimref{std-algebra.witt-basis-car} for each Fock system), so the twisted slot and pair Majoranas are a group of signed directions of order $2^{11}$, which the check confirms directly: closure of order $2048$, the $90$ ordered pairs of distinct directions anticommuting, the squares alternating, the $1023$ nonempty blades traceless, the $1024$ blades distinct up to sign, and the recovered $\tfrac12(\gamma \pm \gamma')$ equal to the Fock operators and satisfying the $50$ relations of the CAR. The two sign bits of the slot and pair systems are both the operator $-1$ of $V$, so sharing one sign bit is automatic on $V$; what (F2) adds to the two separate systems is the anticommutation of the $24$ slot-pair pairs.

\emph{(ii).} (F2) leaves the signature free, but a mode determines it: $\alpha^{2} = 0$ and $\{\alpha, \alpha^{*}\} = 1$ give $\gamma^{2} = 1$ and $\gamma'^{2} = -1$, so five modes give the signature $(5,5)$ and nothing in the clause chooses it. The check computes the squares of all $1024$ blades as operators, finds them equal to GA's blade squares for the sign vector $(1, -1, \dots)$ (the quadratic form $q$ of Theorem \claimref{synthesis.blade-square-signs-f2-quadratic-form}), and the Gauss sum $\sum_{S}\gamma_{S}^{2} = 32 = 2^{5}$, the value of Theorem \claimref{synthesis.signed-directions-and-the-integer-clifford-algebra}(viii) for $p - q \equiv 0, 2 \pmod 8$, shared with $(10,0)$ (also $32$) and not with $(0,10)$ ($-32$); so $\Cl_{\Z}(5,5) \cong \Cl_{\Z}(10,0)$ by (viii), the class of the $\Cl(10)$ of clause (f) (whose generators over $\mathrm{span}\{1, I_{4}\}$ are the five $\gamma$'s and the five $I_{4}\gamma'$'s, Theorem \claimref{synthesis.joint-clifford-clause-is-a-klein-convention-on-the-composite-space}(ii)).

\emph{(iii).} $\gamma^{*} = \gamma = \gamma^{-1}$ and $\gamma'^{*} = \alpha^{*} - \alpha = -\gamma' = \gamma'^{-1}$; the adjoint and the inversion are both anti-automorphisms of $G$ agreeing on generators, so $g^{*} = g^{-1}$ on all of $G$ (the check confirms it on all $2048$ elements, as transposes of signed permutation matrices). The inversion $g \mapsto g^{-1}$ of $\Z[G]$ fixes $c$, hence preserves $(c + 1)$ and descends to an anti-involution of $A$: the adjoint on the image of $A$ is the group's own inversion. A symmetric bilinear form $B$ on the rational span of the occupation basis is $G$-invariant ($g^{\mathsf T}Bg = B$) exactly when $B$ commutes with $G$, since $g^{\mathsf T} = g^{-1}$; the check computes the commutant of the ten Majoranas in $M_{32}(\Q)$ by a signed union-find on the entries ($K_{\pi(i)\pi(j)} = s_{i}s_{j}K_{ij}$ for each signed permutation) and finds it one-dimensional (two-dimensional for nine of them, the control). So the $G$-invariant form is the Fock form up to a scalar, and the adjoint used in the CAR is the one the group fixes.

\emph{(iv).} The pair parity $P = (-1)^{k}$ commutes with the slot operators and anticommutes with the pair operators, and for each mode $\gamma\gamma' = (\alpha + \alpha^{*})(\alpha - \alpha^{*}) = \alpha^{*}\alpha - \alpha\alpha^{*} = -(-1)^{n}$, so $P = -\gamma_{5}\gamma_{6}\cdots\gamma_{10}$ (the six pair Majoranas) lies in $G$. The untwisted slot Majoranas are $\gamma_{i}P$ ($i \le 4$): they lie in $G$, have the same squares ($(\gamma_{i}P)^{2} = \gamma_{i}^{2}P^{2} = \gamma_{i}^{2}$) and commute with every pair Majorana. So the untwisted family generates a subgroup of $G$ of order $2^{11}$ (the check computes it), which is $G$ itself, but it is not a family of signed directions: the $24$ slot-pair pairs commute and (F2c) fails. The group, its sign bit, its signature and its inversion are therefore the same in both conventions, and the content of clause (f) beyond (F2) and (F3) is the choice of which ten elements of $G$ are the directions; by Theorem \claimref{synthesis.joint-clifford-clause-is-a-klein-convention-on-the-composite-space}(iv) that choice is canonical on $V$ up to a central unit. This is the extraspecial central product $2^{1+4}\circ 2^{1+6} \cong 2^{1+10}$ of the slot and pair groups, read with two different generating sets (prior art; not used).

\emph{Status.} Transport-absent. (F1) enters only as the index set $\mathbb{F}_{2}^{10} = G/\langle c\rangle$ of the ten direction labels (not the $2^{5}$ occupation configurations of $V$); identifying these labels with the foundation's distinctions is a reading. The equivalence of the CAR with the Clifford relations of the Majoranas in (i) is classical (Jordan and Wigner), derived here for completeness. Nothing here derives clause (f) from tower structure; it shows that, on $V$, the clause is the statement of Premises (F2) and (F3) for the ten slot and pair Majorana directions, with no further content.
